\chapter{Riesgo y Estructura de Capital}\label{ch:risk-and-capital-structure}

% Alcance: CAPM, beta, riesgo sistemático vs. idiosincrático, Modigliani–Miller, WACC.
% Referencia cruzada: WACC es la tasa de descuento r en el DCF (\cref{ch:valuation}) y en la fórmula de CLV
%   (\cref{ch:customer-economics}).

\section{Riesgo Sistemático y CAPM}\label{sec:capm}
% complexity: 4 -> figure capm-sml

¿Qué retorno debe generar un negocio para justificar el riesgo que soportan sus inversores? CAPM responde
descomponiendo el riesgo en la parte que el mercado valora y la parte que ignora, y derivando luego el
retorno mínimo requerido en consecuencia. Los inversores pueden eliminar el riesgo idiosincrático
(específico de la empresa) de forma gratuita mediante la diversificación. No pueden eliminar el riesgo
sistemático --- la co-variación con el mercado amplio. Dado que solo el riesgo sistemático sobrevive en
una cartera diversificada, solo el riesgo sistemático exige una prima de retorno. La beta mide esa
co-variación; la línea del mercado de valores la valora.

Sea \(R_f\) la tasa libre de riesgo, \(E(R_m)\) el retorno esperado del mercado y \(\beta\)
el riesgo sistemático del activo --- la pendiente de la regresión de sus retornos en exceso sobre los
retornos en exceso del mercado. El modelo de valoración de activos de capital (CAPM) proporciona el
retorno requerido:
\begin{equation}\label{eq:capm}
  E(R_i) \;=\; R_f \;+\; \beta\,\bigl[E(R_m)-R_f\bigr].
\end{equation}
El término \([E(R_m)-R_f]\) es la prima de riesgo de capital (ERP). Para el mercado estadounidense en
horizontes largos, la ERP es aproximadamente \qty{5}{\percent} (media geométrica). Todo activo con
\(\beta>1\) amplifica los movimientos del mercado; \(\beta<1\) los amortigua.

CAPM se basa en preferencias de media-varianza y distribuciones normales de retornos. Los retornos
empíricos tienen colas gruesas; Fama--French documentan que el tamaño y el ratio precio-valor contable
explican variaciones de retorno que \(\beta\) no captura. Las empresas SaaS de alto crecimiento suelen
presentar beneficios contables negativos, lo que hace que las estimaciones históricas de beta sean
ruidosas (los beneficios negativos producen series de retornos poco fiables). Los profesionales suelen
utilizar la beta \emph{sectorial} --- desapalancar la beta de las empresas comparables y re-apalancarla
en la estructura de capital objetivo --- para sortear el ruido específico de cada empresa.

\begin{example}[Derivación del costo del capital propio para la SaaS de referencia]
La SaaS B2B de referencia tiene una \(\beta=1.4\) estimada (superior al mercado, coherente con la
volatilidad propia de las empresas de software en fase de crecimiento). La tasa libre de riesgo actual
(bono del Tesoro a 10 años) es \qty{4}{\percent}; la ERP es \qty{5}{\percent}. Aplicando \cref{eq:capm}:
\[
  E(R_i) \;=\; \qty{4}{\percent} \;+\; 1.4\times\qty{5}{\percent}
          \;=\; \qty{4}{\percent} + \qty{7}{\percent}
          \;=\; \qty{11}{\percent}.
\]
Este \qty{11}{\percent} es el costo del capital propio que alimenta el WACC (\cref{eq:wacc})
y cierra el ciclo de la tasa de descuento del libro: es la \(r\) de \cref{eq:gordon},
\cref{eq:pv}, \cref{eq:npv} y \cref{eq:dcf} en todo el ejemplo de referencia.
Implicación gerencial: elevar las proyecciones de crecimiento de forma que aumente la volatilidad
percibida --- subiendo \(\beta\) de 1.4 a 1.7 --- añade \qty{1.5}{\percent} al retorno requerido,
ensanchando el denominador \((\mathrm{WACC}-g)\) de 0.08 a 0.095
y comprimiendo el múltiplo EV/Ingresos de aproximadamente \(2.6\times\) a \(2.2\times\)
(\cref{eq:ev-revenue}).
\end{example}

\begin{figure}[htbp]
  \centering
  \includegraphics[width=0.86\linewidth]{capm-sml}
  \caption{La Línea del Mercado de Valores. El retorno esperado aumenta linealmente con \(\beta\);
  los activos por encima de la línea están infravalorados (alfa positivo), los que están por debajo
  están sobrevalorados. La SaaS de referencia (\(\beta=1.4\)) se sitúa en el retorno requerido
  del \qty{11}{\percent} que ancla todo el trabajo de DCF en este libro.}
  \label{fig:capm-sml}
\end{figure}

\begin{admonition}{Advertencia}
Tratar \(\beta\) como una propiedad estable e intrínseca. \emph{La beta se estima a partir de la
co-variación histórica}, que cambia con el modelo de negocio, el apalancamiento y el régimen del
mercado. Una empresa SaaS que era defensiva en 2019 (\(\beta=0.8\)) puede revalorizarse hasta
\(\beta=1.6\) cuando las tasas de interés suben y su prima de crecimiento se convierte en el
factor de riesgo dominante. Actualizar la beta en cada ciclo de asignación de capital --- en lugar
de anclarse a una estimación histórica obsoleta --- es la disciplina que mantiene honesta la tasa de
referencia.
\end{admonition}

Los modelos multifactoriales (tres factores de Fama--French, cuatro factores de Carhart, factor de
calidad de AQR) añaden exposiciones de tamaño, valor y momentum al único \(\beta\) del CAPM. Para
empresas privadas sin retornos cotizados, los analistas construyen una beta sintética a partir del
apalancamiento operativo (razón de costos fijos sobre costos totales) y del apalancamiento financiero
--- ambos multiplican directamente la componente sistemática de la volatilidad de los beneficios.

CAPM proporciona el insumo del costo de \emph{equity} que necesita el WACC; los dos son inseparables.
\textcite{brealey2020} derivan el CAPM a partir de la teoría de carteras de media-varianza;
\textcite{ross2022} cubren la literatura empírica y las alternativas factoriales;
\textcite{berk6thcorporate} es la referencia estándar para la metodología de beta sectorial
utilizada en la práctica.

\section{El Costo Promedio Ponderado de Capital}\label{sec:wacc}
% complexity: 3

La empresa está financiada tanto con equity como con deuda. Cada uno tiene un costo y un tratamiento
fiscal diferente. El WACC es la tasa de descuento única que los pondera correctamente para su uso en un
DCF, una regla de NPV o un modelo de CLV. Los accionistas soportan el riesgo residual y exigen un
retorno en consecuencia; los tenedores de deuda tienen prioridad de cobro y aceptan un retorno menor,
adicionalmente subsidiado por la deducibilidad fiscal de los intereses. El WACC combina estos costos
ponderados por su participación en el capital total, produciendo el verdadero costo integral de financiar
los activos de la empresa.

Sean \(E\) y \(D\) los valores de mercado del equity y la deuda, \(V=E+D\), \(r_e\) el costo
del capital propio derivado de \cref{eq:capm}, \(r_d\) el costo de la deuda antes de impuestos y
\(T_c\) la tasa del impuesto corporativo. El escudo fiscal sobre los pagos de intereses reduce el
costo efectivo de la deuda de \(r_d\) a \(r_d(1-T_c)\):
\begin{equation}\label{eq:wacc}
  \mathrm{WACC} \;=\; \frac{E}{V}\,r_e \;+\; \frac{D}{V}\,r_d\,(1-T_c).
\end{equation}
Con \(r_e\) de \cref{eq:capm}, el WACC se convierte en la tasa de descuento unificada que fluye hacia
\cref{eq:dcf}, \cref{eq:terminal-value} y \cref{eq:ev-revenue} en \cref{ch:valuation}.
También equivale a \(r\) en la fórmula de CLV \cref{eq:clv} de \cref{ch:customer-economics} y
en la perpetuidad creciente \cref{eq:gordon}: cada cálculo de valor presente en este libro
utiliza el mismo número subyacente.

El WACC es estable solo si la estructura de capital es estable. Una empresa que ejecuta una
recapitalización apalancada, un LBO o una amortización gradual de la deuda tiene un WACC que varía en
el tiempo; el método del valor presente ajustado (APV) --- que descuenta por separado el FCF no
apalancado y el valor presente del escudo fiscal --- maneja esto con mayor limpieza. Además, el WACC
trata el escudo fiscal como cierto, lo que exige que la empresa sea suficientemente rentable para
aprovecharlo. Una SaaS en fase de crecimiento con EBITDA cercano a cero obtiene poco beneficio del
escudo fiscal, de modo que los primeros años financiados con equity soportan el costo íntegro del
capital propio.

\begin{example}[WACC con y sin deuda]
La SaaS financiada únicamente con equity (\(D=0\)):
\[
  \mathrm{WACC} \;=\; \frac{E}{V}\times\qty{11}{\percent} \;+\; 0
               \;=\; \qty{11}{\percent}.
\]
Supóngase ahora que los fundadores contemplan captar \money{1000000} de deuda al \qty{7}{\percent}
con una tasa impositiva del \qty{25}{\percent}, mientras la empresa se valora en \money{18000000}.
La nueva estructura de capital: \(E=\money{17000000}\), \(D=\money{1000000}\),
\(V=\money{18000000}\):
\[
  \mathrm{WACC} \;=\; \frac{17}{18}\times\qty{11}{\percent}
               \;+\; \frac{1}{18}\times\qty{7}{\percent}\times(1-0.25)
               \;\approx\; \qty{10.68}{\percent}.
\]
La reducción es real pero pequeña --- \qty{32}{\percent} de un punto porcentual ---
porque la deuda es modesta en relación con el valor del equity. Modigliani–Miller
(\cref{sec:mm}) formaliza por qué esta cuasi-irrelevancia se sostiene en su forma más pura,
y qué la rompe en la práctica.
\end{example}

\begin{admonition}{Advertencia}
Usar una estructura de capital a \emph{valor contable} en lugar de \emph{valor de mercado} en
\cref{eq:wacc}. Los valores contables reflejan la contabilidad histórica, no el costo de oportunidad
del capital actual. Una empresa con \money{5000000} de equity en libros pero una capitalización de
mercado de \money{50000000} está financiada predominantemente con equity a pesos de mercado;
aplicar pesos contables sobreestima drásticamente la fracción de equity y sobreestima el WACC,
llevando a la empresa a rechazar proyectos con NPV positivo.
\end{admonition}

Para empresas con estructuras de capital complejas (múltiples tramos de deuda, equity preferente,
garantías), cada instrumento entra en \cref{eq:wacc} con su propio costo y peso de mercado.
Las notas convertibles requieren separar los componentes de deuda y equity (valor Black--Scholes
de la opción de compra incorporada) antes de ponderarlos. La estimación del WACC para empresas
privadas añade otra capa: dado que ni el equity ni la deuda cotizan en bolsa, tanto el costo del
capital propio (utilizando el método de beta sectorial) como el costo de la deuda (utilizando una
calificación sintética a partir de razones de cobertura de intereses) deben estimarse en lugar de
observarse.

El WACC es la columna vertebral de la arquitectura de tasas de descuento del libro: su componente de
equity proviene de \cref{eq:capm}; alimenta \cref{eq:dcf} y \cref{eq:terminal-value} en
\cref{ch:valuation}; y equivale a \(r\) en el modelo de economía del cliente
(\cref{eq:clv}, \cref{eq:gordon}).
\textcite{brealey2020} derivan el ajuste por escudo fiscal; \textcite{modigliani1958capital}
es la proposición original.

\section{Modigliani--Miller}\label{sec:mm}
% complexity: 4

¿Importa la forma en que se financia la empresa --- equity versus deuda? Modigliani–Miller establece
las condiciones bajo las cuales no importa y, al mismo tiempo, las imperfecciones que hacen que la
estructura de capital sí sea relevante. Comprender el teorema es un prerrequisito para saber cuándo
ignorarlo. En un mundo sin impuestos, sin costos de quiebra, con información perfecta y sin fricciones
de agencia, el valor total de la empresa está determinado únicamente por sus activos reales --- no por
cómo divide las reclamaciones sobre esos activos en deuda y equity. La financiación es entonces una
redistribución de suma cero, no un acto de creación de valor.

\textbf{Proposición I (sin impuestos):} \(V_L = V_U\) --- el valor de una empresa apalancada
equivale al de una empresa idéntica no apalancada. \textbf{Proposición II (sin impuestos):} a medida
que el apalancamiento aumenta, los accionistas exigen un retorno mayor para compensar el riesgo
adicional que asumen:
\[
  r_e \;=\; r_A \;+\; \frac{D}{E}(r_A - r_d),
\]
donde \(r_A\) es el costo del capital (de los activos) no apalancado. Con impuestos corporativos, el
escudo fiscal sobre intereses tiene un valor presente de \(T_c\cdot D\), que incrementa el valor de
la empresa:
\begin{equation}\label{eq:mm-wacc}
  V_L \;=\; V_U \;+\; T_c\,D,
  \qquad
  \mathrm{WACC} \;=\; r_A\,\bigl(1 - T_c\,D/V\bigr).
\end{equation}
\cref{eq:mm-wacc} muestra el WACC disminuyendo monotónicamente con el apalancamiento cuando solo se
consideran los impuestos --- lo que implicaría un valor máximo al \qty{100}{\percent} de deuda, lo cual
es evidentemente incorrecto. Las piezas que faltan son el objeto de \cref{sec:optimal-capital-structure}.

El mundo sin fricciones de Modigliani–Miller omite tres fuerzas que dominan con apalancamiento
moderado a elevado:
(1) \emph{costos de dificultad financiera} --- costos legales, operativos y reputacionales
desencadenados cuando la deuda se vuelve difícil de servir; (2) \emph{costos de agencia} --- a medida
que la deuda aumenta, los accionistas obtienen un incentivo similar al de una opción para asumir riesgos
excesivos (\textcite{jensen1976agency} lo formalizan como el problema de sustitución de activos);
(3) \emph{información asimétrica} --- la teoría de la jerarquía de financiación
(\textcite{brealey2020}) predice que las empresas prefieren primero los beneficios retenidos, luego la
deuda y por último el equity, porque emitir equity señala que la dirección cree que las acciones
están sobrevaloradas. Cada fricción crea una brecha entre la irrelevancia de Modigliani–Miller y el
comportamiento financiero observado.

\begin{example}[Escudo fiscal de Modigliani–Miller para la SaaS]
Con la estructura de capital observada de la SaaS (\(D=0\), \(V=\money{18000000}\)) y una
tasa impositiva del \qty{25}{\percent}, \cref{eq:mm-wacc} da:
\[
  \mathrm{WACC} \;=\; \qty{11}{\percent}\times(1-0.25\times0) \;=\; \qty{11}{\percent}.
\]
El escudo fiscal es cero porque no hay deuda. Si la empresa añadiera
\money{5000000} de deuda permanente:
\[
  \mathrm{WACC} \;\approx\; \qty{11}{\percent}\times\!\left(1-0.25\times\tfrac{5}{18}\right)
               \;\approx\; \qty{10.24}{\percent}.
\]
El WACC menor aumenta el valor de la empresa de forma mecánica, pero solo hasta el punto en que
los costos de dificultad financiera comienzan a compensar el escudo fiscal --- la compensación que
\cref{sec:optimal-capital-structure} desarrolla.
\end{example}

\begin{admonition}{Advertencia}
Concluir a partir de Modigliani–Miller que la estructura de capital \emph{nunca} importa. El valor
del teorema no radica en su conclusión en el mundo sin fricciones, sino en su enumeración de las
fricciones que hacen que la estructura de capital sí sea relevante en el mundo real: impuestos,
costos de dificultad financiera, costos de agencia e información asimétrica. \emph{Toda decisión
real de estructura de capital es una apuesta sobre qué fricción domina.}
\end{admonition}

El marco de costos de agencia de Jensen--Meckling añade dos distorsiones adicionales más allá de
\textcite{modigliani1958capital}: la subinversión (el sobreendeudamiento reduce el incentivo del
accionista para financiar proyectos con NPV positivo, porque las ganancias se acumulan en favor de los
bonistas) y la sobreinversión/desplazamiento de riesgo (los accionistas prefieren proyectos arriesgados
porque su posición se asemeja a una opción de compra con pérdidas limitadas). Estos efectos crean un
nivel \emph{óptimo} de deuda incluso en ausencia de impuestos: aquel que minimiza los costos totales
de agencia.

Modigliani–Miller sigue siendo el fundamento de la teoría de estructura de capital precisamente porque
obliga a esclarecer \emph{qué} fricciones importan en cada situación. \textcite{modigliani1958capital}
es el trabajo original; \textcite{jensen1976agency} extiende el análisis a los costos de agencia;
\textcite{brealey2020} proporciona el tratamiento unificado para profesionales.

\section{Estructura de Capital Óptima}\label{sec:optimal-capital-structure}
% complexity: 5 -> figure capital-structure

¿Cuánta deuda debería asumir la empresa? La teoría del equilibrio sostiene que la respuesta está donde
el beneficio marginal del escudo fiscal iguala el costo marginal de la dificultad financiera. La teoría
de la jerarquía de financiación sostiene que la pregunta es casi irrelevante: las empresas se financian
secuencialmente a partir del efectivo interno primero. Agregar deuda crea dos fuerzas opuestas: el
escudo fiscal eleva el valor de la empresa; los costos de dificultad financiera lo reducen. La razón de
apalancamiento que maximiza el valor se sitúa en el pico de la curva de beneficio neto --- donde el
siguiente dólar de deuda añade más costo de dificultad financiera de lo que ahorra en impuestos. Para
una empresa rentable y estable, ese pico puede estar en un apalancamiento sustancial; para una empresa
en crecimiento con pérdidas, está cerca de cero.

La teoría del equilibrio define la estructura de capital óptima como:
\[
  V^* \;=\; V_U \;+\; PV(\text{escudo fiscal}) \;-\; PV(\text{costos de dificultad financiera}).
\]
Este es un marco cualitativo: ni la función de costos de dificultad financiera ni la función de costos
de agencia tiene solución de forma cerrada, de modo que la «derivación» es conceptual. \cref{eq:mm-wacc}
muestra el escudo fiscal aumentando monotónicamente con \(D/V\); los costos de dificultad financiera
son convexos en el apalancamiento (pequeños con deuda baja, crecientes abruptamente una vez que las
razones de cobertura se deterioran). El óptimo es un nivel de \(D/V\) que no puede mejorarse mediante
un intercambio marginal de deuda por equity.

La teoría del equilibrio asume que la empresa puede predecir su función de costos de dificultad
financiera, que depende de la tangibilidad de los activos (los activos físicos son más fáciles de
liquidar, de modo que los costos son menores), la estabilidad de los beneficios y el costo de la
renegociación. Los negocios intangibles de alto crecimiento (SaaS, biotecnología) tienen costos de
dificultad financiera elevados porque sus activos (relaciones con clientes, propiedad intelectual,
personal clave) se evaporan en situaciones de dificultad --- lo que empuja el \(D/V\) óptimo cerca
de cero. La teoría de la jerarquía de financiación ofrece una explicación alternativa: las empresas no
optimizan \(D/V\); simplemente utilizan la financiación más barata disponible, lo que para las
empresas rentables significa primero el efectivo interno.

\begin{example}[Estructura óptima para la SaaS en fase de crecimiento]
La SaaS en fase de crecimiento tiene EBITDA negativo en los años 1--2, sin ingresos gravables que
proteger, y activos de relación con clientes que perderían valor rápidamente en una situación de
dificultad financiera (el \emph{churn} se acelera; el talento se va). La teoría del equilibrio predice
por tanto un \(D/E\approx0\) óptimo, que es precisamente la financiación observada: financiada con
equity a través de capital de riesgo, sin deuda. El WACC del \qty{11}{\percent} es el costo puro del
capital propio (\cref{eq:capm}).
La deuda pasa a ser relevante solo cuando la empresa alcanza una rentabilidad estable y obtiene un
escudo fiscal que vale la pena proteger. Implicación gerencial: la decisión de mantenerse sin
apalancamiento no es timidez --- es la lectura correcta del equilibrio para un negocio cuyo valor
está concentrado en intangibles y opciones de crecimiento.
\end{example}

\begin{figure}[htbp]
  \centering
  \includegraphics[width=0.86\linewidth]{capital-structure}
  \caption{Valor de la empresa frente a apalancamiento según la teoría del equilibrio.
  El escudo fiscal eleva el valor por encima de la línea de referencia no apalancada;
  los costos de dificultad financiera (convexos en el apalancamiento) lo erosionan con
  un \(D/V\) elevado.
  El pico es la estructura de capital óptima. Una SaaS en fase de crecimiento
  con activos intangibles y beneficios volátiles se sitúa cerca de \(D/V=0\).}
  \label{fig:capital-structure}
\end{figure}

\begin{admonition}{Advertencia}
Comparar la estructura de capital con los promedios sectoriales sin ajustar por la etapa y la
composición de activos propios de la empresa. Una empresa de software madura con márgenes brutos
del \qty{60}{\percent} y flujo de caja libre estable puede soportar un apalancamiento significativo;
una SaaS en fase previa a la rentabilidad no puede, aunque opere en el mismo sector. \emph{Las
razones de apalancamiento sectoriales son una descripción, no una prescripción.}
\end{admonition}

La teoría del \emph{market timing} (\textcite{brealey2020}) propone que los gestores emiten equity
cuando las valoraciones son altas y recompran cuando son bajas --- sin prestar atención a una razón
objetivo óptima. La evidencia empírica de que el apalancamiento revierte a la media lentamente (en
lugar de hacerlo rápidamente hacia un objetivo) es coherente tanto con el \emph{market timing} como
con los costos de ajuste. Los modelos dinámicos de equilibrio permiten que el \(D/V\) objetivo varíe
a medida que la empresa pasa de la fase de crecimiento a la madurez, capturando el patrón del ciclo
de vida: equity de capital de riesgo → capacidad de deuda rentable → recapitalización apalancada →
recompra de acciones.

La estructura de capital es el puente entre el marco de costo de capital desarrollado en
\cref{sec:capm,sec:wacc} y la mecánica de valoración de \cref{ch:valuation}: la estructura de
capital correcta minimiza el WACC y maximiza el valor de la empresa, que es lo que el DCF de
\cref{eq:dcf} y los múltiplos de \cref{eq:ev-revenue} están midiendo.
\textcite{brealey2020} examinan la literatura empírica sobre la elección de la estructura de capital;
\textcite{jensen1976agency} desarrolla la justificación de los costos de agencia para un óptimo
interior.
